\section{任务完成超过评分覆盖}
\label{sec:evaluation}

\textbf{实验协议。} 三个云端端点在同一个智能体修订版下处理相同的
\CorpusTaskCount{} 个均衡任务套件。每个完成输出必须通过 \CSim、\Synth、
100\,MHz、U55C 容量以及 Track-A 契约~\cite{fpt26tracka} 下所有适用的
\CoSim{} 门。模型覆盖三种架构—规模体制：DeepSeek V4 Pro（1.6\,T MoE，
49\,B 激活）、Qwen3.5-122B-A10B（122\,B MoE，10\,B 激活）和 Qwen3.6-27B
（27\,B 稠密）。

\begin{table}[!h]
  \centering
  \caption{\textbf{最小的稠密模型领先。} QoR 分数估计为
  $100\,Q_{\mathrm{HW}}\,E$；DeepSeek 未返回可测量的 $Q_{\mathrm{HW}}$ 提升。}
  \label{tab:headline-results}
  \scriptsize
  \resizebox{\columnwidth}{!}{%
  \begin{tabular}{lrrrrr}
    \toprule
    端点 & 完成数 & 令牌（M） & 积分 & 估计 QoR 均值 & $>$76 比例 \\
    \midrule
    \MainModelResultsRows
    \bottomrule
  \end{tabular}}
\end{table}

\textbf{结果。} 表~\ref{tab:headline-results} 总结了三次实验。
Qwen3.6-27B 取得了 \QwenThreeSixSuccessRate{} 的完成率
（\QwenThreeSixSuccessCount{}/\CorpusTaskCount{} 个任务，
\QwenThreeSixFailureCount{} 次失败），包括在最具挑战性的结构 \CoSim{} 修复
类别中取得 25/25 满分。尽管是最小的模型，它在完成率和优化质量上均超越了
两个更大的 MoE 模型。Qwen3.5-122B-A10B 在令牌效率上最优
（\QwenThreeFiveTokensM{}M 令牌），而 DeepSeek V4 Pro 消耗了
\DeepSeekTokensM{}M 令牌（是 Qwen3.6 的 3.1$\times$），却未产生可测量的
QoR 提升。

逐类别分解（附录~\ref{app:results}）揭示四种模式。
（1）仅显式 QoR 优化任务产生了超越基线的分数；全部五个修复/生成类别均聚集在
$Q_{\mathrm{HW}}{\approx}0.7576$（$S{\approx}73.7$，$E{=}0.973$）附近。
智能体通过了验证门，但未在修复任务上执行架构级优化。
（2）Qwen3.6-27B 在优化质量上占优（平均 $Q_{\mathrm{HW}}{=}0.812$，60\%
的优化任务 $>$76，52\% $>$80），而 Qwen3.5 平均为 $Q_{\mathrm{HW}}{=}0.775$
（28\% $>$76）。
（3）对 HLS 修复任务而言，模型架构似乎比参数量更具预测力：27\,B 稠密模型
在结构修复上超越了 122\,B 和 1.6\,T 两个 MoE 模型（100\% vs.\ 80\% 和
92\%）。

\takeaway{Qwen3.6-27B 在完成率（\QwenThreeSixSuccessRate）和优化质量
（QoR 任务平均 \QwenThreeSixQorScore{} 分）上均领先，表明对 HLS 修复而言
架构比参数量更重要。仅显式优化任务受益于 LLM 驱动的 QoR 提升；修复任务
达到了正确性但未获得质量收益。}
